\section{Methodology}

\subsection{Environment}

Each game uses four agents, a standard 52-card deck, and a seeded shuffle. The player holding the Ace of Spades moves first. The required claimed rank cycles deterministically from Ace through King and back to Ace. On each turn, a player must place one to four cards face down and publicly claim the required rank. Players may lie about rank but not about card count. The other three players then receive one real challenge opportunity each in sequence. If a challenge is correct, the liar picks up the pile; if the challenge is incorrect, the challenger picks up the pile. The first player to empty their hand wins.

\subsection{Prompt Conditions}

The pilot uses four conditions. Experiment~0 is a low-strategy reference rather than a true random baseline. Experiment~1 explicitly legalizes and normalizes deception. Experiment~2 tells the focal model it may lie while opponents are honesty-constrained, providing an asymmetric-fairness framing. Experiment~3 is a plain-language honesty mandate that explicitly forbids lying. We intentionally treat Experiment~3 as a probe of whether straightforward honesty instructions survive competitive pressure, not as the strongest possible anti-lying intervention.

\subsection{Model Cohort}

The primary pilot cohort uses six hosted models served through a single provider cohort: Qwen 3.5 397B A17B, MiniMax M2.5, Nemotron 3 Super 120B A12B, Mistral Small 4 119B, GLM5, and Kimi K2.5. The full pilot design spans all $\binom{6}{4}=15$ unique four-model matchups, with ten seeded games per matchup per experiment for a target of 600 games.

\subsection{Implementation and Dataset Hygiene}

The environment is implemented in TypeScript with structured-response parsing, bounded prompt context, provider retry logic, client recreation on recoverable failures, and long-run recovery for hosted inference. The public release also exposes an importable environment API with \texttt{reset}, \texttt{observation}, \texttt{publicState}, \texttt{step}, \texttt{done}, and \texttt{result} methods, plus a small baseline-policy pack for local side comparisons. Reproducibility and dataset hygiene are enforced through seeded deck order and seating, prompt versioning and hashing, schema versioning, dominant-cohort filtering at analysis time, and versioned dataset/evaluation manifests with checksums. Research-default runs are uncapped and must end in a natural win; capped runs are marked explicitly and excluded from the main cohort.

\subsection{Analysis}

The primary statistical unit is one row per player per game rather than one row per turn. This supports more defensible summaries of win rate, lie frequency, lie success, challenge frequency, challenge accuracy, and instruction-violation rate. Qualitative reasoning traces and exemplar games are used as supporting evidence rather than as the main statistical unit.

\begin{figure}[t]
  \centering
  \includegraphics[width=\linewidth]{benchmark_overview.png}
  \caption{Bullshit-Bench overview. Four LLM agents play seeded games of \emph{Bullshit} under one of four prompt conditions. The engine records the actual cards played, the public claim, whether the move was a lie, whether challenge opportunities were offered, and the outcome of any challenge. These provenance-aware logs are then aggregated into player-game metrics, frozen cohort artifacts, and paper figures.}
  \label{fig:benchmark-overview}
\end{figure}
